% Generated draft: requires mathematical and source review.
\documentclass[11pt]{article}
\usepackage[margin=25mm]{geometry}
\usepackage{fontspec}
\setmainfont{Latin Modern Roman}
\usepackage{amsmath,amssymb,mathtools,unicode-math}
\setmathfont{Latin Modern Math}
\usepackage{enumitem,needspace,xurl,hyperref}
\setlength{\parindent}{0pt}
\setlength{\parskip}{5pt}
\setlength{\emergencystretch}{3em}
\newcommand{\N}{\mathbb{N}}
\newcommand{\Z}{\mathbb{Z}}
\newcommand{\Q}{\mathbb{Q}}
\newcommand{\R}{\mathbb{R}}
\newcommand{\C}{\mathbb{C}}
\newcommand{\F}{\mathbb{F}}
\newcommand{\A}{\mathbb{A}}
\newcommand{\PP}{\mathbb{P}}
\newcommand{\E}{\mathbb{E}}
\newcommand{\eps}{\varepsilon}
\newcommand{\dd}{\,\mathrm{d}}
\begin{document}
\section*{UnsolvedMath ID: 6200097}
\textbf{Source problem code:} catalogue entry
\subsection{Definitions and mathematical statement}
A nonempty topological space \(X\) is contractible if there exist \(x_0\in X\) and a continuous homotopy \(H:X\times[0,1]\to X\) with
\[
H(x,0)=x,\qquad H(x,1)=x_0\quad\text{for all }x\in X.
\]
Ilia Kapovich's question asks for a uniform terminating decision algorithm which, given a finite cell complex in an effective finite representation, outputs yes exactly when its realization is contractible. No bound on dimension or additional promise such as simple connectivity is stated.

A completely finite input subclass is that of finite abstract simplicial complexes, encoded by their vertex set and maximal faces. The corresponding explicit decision target is a Turing algorithm \(D\) satisfying, for every such complex \(K\),
\[
D(K)\text{ halts},\qquad D(K)=1\iff |K|\text{ is contractible}.
\]
Any algorithm for the full stated cell-complex class must handle this subclass. For a general finite CW description the attaching maps also need a finite effective encoding; arbitrary continuous attaching maps are not themselves finite algorithmic inputs. A search that halts only when a contraction is found does not meet the requested decision requirement.

The supplied research summary answers the general algorithm question negatively and invokes Adian--Rabin/homotopy undecidability, while the stored status remains open and OPEN-TRIAGE. These contradictory metadata are retained explicitly. The summary gives no precise contractibility theorem, input category, dimension range, or reduction, so its conclusion is a supplied literature assessment rather than an independently verified undecidability result. Triviality of the fundamental group alone is not substituted for contractibility.
\subsection{Short English statement}
Is contractibility of finite cell complexes decidable by an algorithm that halts on every valid finite input? The supplied research summary says no, but the registry retains an open label and provides no precise undecidability citation.
\subsection{Sources}
\begin{itemize}
\item Local UnsolvedMath registry, permanent ID \(6200097\), exact source code \texttt{catalogue entry}, ``Boundaries of Groups and Kleinian Groups --- Problem 97.'' Proposer Ilia Kapovich, 2005.
\item Kapovich, \emph{Problems on Boundaries of Groups and Kleinian Groups} (2005), Problem 97, recorded PDF page 24; \url{https://www.math.ucdavis.edu/~kapovich/EPR/problems.pdf}. Recorded access 2026-07-29; pdf-text-normalized extraction.
\item Supplied research summary mentions Adian--Rabin and homotopy/fundamental-group undecidability but supplies no complete publication reference or exact contractibility theorem.
\end{itemize}
\end{document}
